\documentclass[11pt]{article}
\usepackage[a4paper,margin=1in]{geometry}
\usepackage[T1]{fontenc}
\usepackage{newtxtext,newtxmath}
\usepackage{microtype}
\usepackage{amsmath}
\usepackage{booktabs,tabularx}
\usepackage[font=small,labelfont=bf]{caption}
\usepackage{enumitem}
\usepackage{xcolor}
\usepackage[colorlinks=true,linkcolor=black!70!blue,citecolor=black!70!blue,urlcolor=black!70!blue]{hyperref}
\usepackage[capitalise,noabbrev]{cleveref}

\setlist{itemsep=2pt,topsep=4pt}
\newcommand{\ip}[2]{\langle #1,#2\rangle}
\newcommand{\norm}[1]{\lVert #1\rVert}
\newcommand{\lean}[1]{\texttt{\small #1}}
\newcommand{\ms}{\,\mathrm{m/s}}
\newcolumntype{Y}{>{\raggedright\arraybackslash}X}

\title{\textbf{Buy Headroom, Not Cleverness}\\[4pt]
\large An Engineering Guide to Quadrotor Pursuit of Fast Fixed-Wing Aircraft}
\author{Abdullah Almanei}
\date{Draft, 2 November 2026}

\begin{document}
\maketitle

\begin{abstract}
A quadrotor must keep a fast fixed-wing aircraft in the image of one camera in each frame.
The aircraft is 3.5\,m long, flies at speeds up to $55\ms$ and turns on a radius of $350$\,m.
A computer has checked all the bounds that we use.
One bound is sufficient: if it is true, our pursuit method keeps the target in view.
The other bound is necessary: if it is false, no method can keep the target in view.
We changed each design parameter and recorded the effect on the two bounds.

The results give three rules.
Thrust and air set the top speed, and the thrust-to-weight ratio is the most effective parameter.
Yaw authority sets the agility in a dive at mid speed, and more thrust does not help there.
The camera position solves the camera problem: fly below the target and point a hinged camera up and back.
More complex pursuit methods did not give better results than a simple copy of the target path.

For a 1.5\,kg quadrotor, the proved limit is $57.7\ms$, and our proofs certify our method to $52\ms$.
A recommended set of changes increases these values to $73.9\ms$ and $68\ms$.
\end{abstract}

\section{Introduction}

A quadrotor that follows a fast fixed-wing aircraft must keep the aircraft in view at all times.
This looks like a problem of control.
Our results show that it is almost always a problem of hardware.

We made a full rigid-body model of the quadrotor and of the target.
In the Lean~4 proof assistant~\cite{lean4,mathlib}, we proved two inequalities about this model.
The first inequality is sufficient for one pursuit method.
The second inequality is necessary for all pursuit methods.

Then we changed each physical parameter and recorded how the two inequalities moved.
This guide gives the results in a form that an engineer can use.
It tells you which parameters to change, by how much, and which parameters to leave alone.

Our reference quadrotor has a mass of 1.5\,kg and a thrust-to-weight ratio of 6.
It has four motors in a plus layout, on arms of 18\,cm.
Its yaw reaction coefficient is 0.02\,m, and its inertia is $(0.012, 0.012, 0.02)$\,kg\,m$^2$.

Its parasitic drag coefficient $c$ is $0.008\,\mathrm{m}^{-1}$.
The rotor disc also has a normal-force coefficient $k_n$ of $0.01\,\mathrm{m}^{-1}$.
This value agrees with approximately $0.02$\,m$^2$ of normal area.
The tilt limit is $85^\circ$, and the body-rate limit is $20$\,rad/s.

A camera with a $120^\circ$ field of view is on the top of the quadrotor.
The camera points to the rear, on a hinge with a range of $30^\circ$ to $70^\circ$.
The control system operates at $240$\,Hz.

The reference target flies at speeds from $20\ms$ to a maximum speed $v_{\max}$.
Its minimum turn radius is $350$\,m.
Its acceleration along the path is not more than $2$\,m/s$^2$.
Its vertical acceleration is not more than $1.5$\,m/s$^2$.
Its snap is not more than $15$\,m/s$^4$, and its jerk is between $1$ and $5$\,m/s$^3$.
It can climb or dive at angles of up to $5^\circ$ or $10^\circ$.

\Cref{tab:summary} gives the answer on one page.
The quadrotor gets the most speed from more thrust and from a smaller rotor-disc area.
It gets the most agility in steep dives from more yaw authority and less yaw inertia.
A hinge or a wide lens removes the camera as a limit.
A longer arm, a cleaner airframe, wings and a more complex planner do not help.
The subsequent sections explain each of these results.

\begin{table}[t]
\centering
\caption{Summary for the reference quadrotor. The limit is the proved necessary bound. The
certified value is the highest $v_{\max}$ for which the proved sufficient condition is true.}
\label{tab:summary}
\begin{tabularx}{\linewidth}{@{}Yl@{}}
\toprule
Question & Answer \\
\midrule
Highest target speed that any method can follow (proved limit) & $57.7\ms$ \\
Certified speed, dives to $5^\circ$, jerk to $5$ & $52.1\ms$ \\
Certified speed, dives to $10^\circ$, jerk to $3$ & $52.1\ms$ \\
Certified speed, dives to $10^\circ$, jerk $5$ & $33.4\ms$ (yaw limit) \\
\quad with yaw coefficient $0.03$ or $30\%$ less inertia & $52.1\ms$ \\
Recommended set of changes (\cref{sec:checklist}) & limit $73.9$; certified $68.1$ and $58.4$ \\
Limit at top speed & the top motor is at maximum thrust \\
Limit in dives & yaw torque at low total thrust \\
Is the camera a limit? & no, with a $120^\circ$ lens or a hinge \\
\bottomrule
\end{tabularx}
\end{table}

\section{How to read the numbers}

This guide uses two types of bound.
The necessary bound is a speed limit.
If the target flies faster than this limit, a straight and level target escapes from all
quadrotor paths.
This is true also for a path that the quadrotor plans with full knowledge of the target's future.
The Lean theorems are \lean{master\_necessary\_power} for the rotor-disc model and
\lean{master\_necessary3} for the isotropic model.

The sufficient bound is the inequality \lean{MasterIneq3}.
If it is true, our method keeps each permitted target in view in each frame.
The target must start in steady flight.
The Lean theorem is \lean{master\_sufficient3}.
The certified speed is the highest $v_{\max}$ for which \lean{MasterIneq3} is true.
A numerical copy of the Lean definitions finds this value.

The model is a full rigid-body quadrotor in discrete time, with a step of $1/240$\,s.
It includes four motors with limits, the motor mixer and the rotational equation of motion.
It also includes exact rotations and the limits on tilt, body rate and speed.
Air causes two forces.
The first force is parasitic drag, $-c\norm v v$.
The second force is a normal force on the rotor disc~\cite{hoffmann2007,faessler2018}:
\begin{equation}
  a_n = -k_n \norm v \ip{v}{b_3}\, b_3 .
\end{equation}
Here $v$ is the velocity and $b_3$ is the thrust axis.
This force acts at the centre of mass.
The target is in view when it is in the detection range and in the field of view.

The model does not include five effects.
These effects are the change of motor thrust with airspeed, aerodynamic moments, wind, the ground
and sensor noise.
Thus each number in this guide is true for the model only.
\Cref{sec:measure} tells you what to measure before you use these numbers.

\section{Geometry: fly below the target and look back and up}
\label{sec:geometry}

\subsection{Why the quadrotor must fly below the target}

A fast quadrotor tilts a large amount.
The pitch angle is approximately $50^\circ$ at $35\ms$ and $70^\circ$ to $80^\circ$ at $55\ms$.
A camera below the body, with the quadrotor above the target, cannot keep the target in view
above approximately $30\ms$.
At these speeds, the tilt of the thrust axis moves the target out of the field of view.

A position above the target is also a bad position for a different reason.
To look down, the quadrotor must pitch more than its drag balance.
This extra pitch pushes the quadrotor forward.
The quadrotor then passes the target and loses it.
The theorem \lean{QuadTraj.forced\_loss\_time} gives the time until this loss.

In a prototype with a range of 382\,m and a $90^\circ$ lens, the time was approximately 17\,s.
This was at $20\ms$, 300\,m above the target.
At $30\ms$, the time was approximately 53\,s.

Below the target, the opposite occurs.
When the quadrotor pitches its nose down for speed, a rear camera on its top turns up.
Thus, speed helps this camera.

\subsection{Point the camera to the rear}

On a quadrotor with its nose down, a forward camera on the top points at the angle $\eta - p$.
Here $p$ is the pitch angle and $\eta$ is the camera angle.
At high speed, this angle is zero or less, so the camera cannot look up.

A rear camera points at the angle $p + \eta$ from the rear horizon.
At high speed, this camera points straight up.
Install the camera on the rear face of a cone-shaped head.
Attach the camera with a hinge that turns about the pitch axis.

\subsection{Use a hinge or a wide lens}

Our method locks the yaw to the side direction of the target.
With this method, the camera error angle is not more than the bank angle plus the hinge error.
The bank angle is not more than $38^\circ$ in all conditions.
The camera error is not more than $40^\circ$.
\Cref{tab:camera} shows when the hinge is necessary.

\begin{table}[h]
\centering
\caption{Certified $v_{\max}$ in m/s. Isotropic drag, dives to $5^\circ$, jerk to $3$.}
\label{tab:camera}
\begin{tabular}{@{}lcc@{}}
\toprule
Camera installation & $120^\circ$ lens & $90^\circ$ lens \\
\midrule
Hinge, $30^\circ$ to $70^\circ$ & $69.9$ & $69.9$ \\
Hinge, $20^\circ$ to $80^\circ$ & $69.9$ & $69.9$ \\
Fixed at $50^\circ$ & $69.9$ & $\mathbf{50.1}$ \\
Fixed at $60^\circ$ & --- & $\mathbf{34.6}$ \\
\bottomrule
\end{tabular}
\end{table}

With a lens of $120^\circ$ or more, a fixed camera at approximately $50^\circ$ is sufficient.
With a lens of $90^\circ$, use a hinge with a range of $30^\circ$ to $70^\circ$ or more.
The pitch angle changes by approximately $80^\circ$ between $20\ms$ and $55\ms$.
Only a hinge or a wider field of view can compensate for this change.

\subsection{Vertical distance and ground clearance}

The vertical distance between the quadrotor and the target has almost no effect on the bounds.
Distances from 60\,m to 150\,m give the same results.
But this distance sets the ground clearance.
The target must fly higher than this distance plus a safety margin.
Use a distance of 60\,m to 100\,m.

\section{Top speed: thrust and air}
\label{sec:speed}

\subsection{The physics}

At high speed, the rotor disc is almost perpendicular to the airflow.
At $55\ms$, its tilt is approximately $68^\circ$.
Thus, the normal force on the disc is large.
It is approximately equal to the thrust that the quadrotor uses against parasitic drag.

Let $F$ be the total thrust per unit mass, with $F \le F_{\max}$.
Let $y = \ip{v}{b_3}$.
The power that the motors can put into the flight path is then limited:
\begin{equation}
  F y - k_n\norm v y^2 \;\le\; \min\!\Big(F_{\max}\norm v,\; \frac{F_{\max}^2}{4k_n\norm v}\Big).
  \label{eq:power}
\end{equation}
Parasitic drag removes the power $c\norm v^3$.
The quadrotor must also stay in camera range of a level target, so its height cannot change
without limit.
The energy balance then gives the proved speed limit (\lean{no\_traj\_keeps\_of\_power}).
To the first order, this limit is
\begin{equation}
  V_{\mathrm{limit}} \approx \Big(\frac{F_{\max}^2}{4\,k_n\,c}\Big)^{1/4}.
  \label{eq:wall}
\end{equation}

\subsection{Effect of each parameter}

\Cref{tab:speed} shows the limit and the certified speed for each parameter value.
\Cref{eq:wall} agrees closely with the data.
When $F_{\max}$ doubles, the limit increases by a factor of $1.42$.
The equation gives $\sqrt2$, which is $1.41$.
When $k_n$ increases by a factor of four, the limit decreases by a factor of $0.71$.
The equation gives $0.707$.

When $c$ increases from $0.006$ to $0.010$, the limit decreases by a factor of $0.88$.
The equation gives the same value.
The certified speed is approximately $90\%$ of the limit in all cases.

\begin{table}[h]
\centering
\caption{Proved limit and certified $v_{\max}$ in m/s. Dives to $5^\circ$, jerk to $5$. All other
parameters have their reference values.}
\label{tab:speed}
\begin{tabular}{@{}llcc@{}}
\toprule
Parameter & Value & Limit & Certified \\
\midrule
Thrust-to-weight ratio & $4$ & $47.1$ & $37.6$ \\
 & $5$ & $52.6$ & $46.2$ \\
 & $6$ & $57.7$ & $52.1$ \\
 & $7$ & $62.3$ & $57.3$ \\
 & $8$ & $66.7$ & $61.9$ \\
\midrule
$k_n$ (m$^{-1}$) & $0$ & $84.5$ & $69.5$ \\
 & $0.005$ & $68.4$ & $61.6$ \\
 & $0.01$ & $57.7$ & $52.1$ \\
 & $0.02$ & $48.7$ & $40.8$ \\
\midrule
Drag $c$ (m$^{-1}$) & $0.006$ & $62.0$ & $39.3$ \\
 & $0.008$ & $57.7$ & $52.1$ \\
 & $0.010$ & $54.5$ & $50.1$ \\
\bottomrule
\end{tabular}
\end{table}

\subsection{How close the limit is to the physics}

The limit of $57.7\ms$ is a physical limit.
In steady, straight and level flight at maximum thrust, the quadrotor gets to $57.3\ms$.
The difference of $0.4\ms$ comes from the energy gain of the time step at $240$\,Hz.

A point mass can fly faster than steady flight with a weave maneuver.
It applies more thrust than necessary and turns the extra thrust to the left and to the right in
turn.
The side force does no work, but it tilts the rotor disc away from the airflow.
This decreases the drag on the disc.
At $k_n = 0.02$, the weave gives $47.8\ms$, and steady flight gives $47.0\ms$.

A real quadrotor cannot fly this weave.
For the fast changes of bank angle, motor thrusts from $-34$\,N to $+78$\,N are necessary.
The motor limit is 22\,N.
Also, this torque is necessary when the total thrust is already at maximum.
In our search, each weave within the limits was slower than steady flight.

\subsection{Rules for top speed}

Increase the thrust-to-weight ratio first.
It is the only parameter that increases both the limit and the certified speed by approximately
$5\ms$ per unit.

Then decrease the area of the rotor plane that faces the airflow.
Use small rotor discs for the thrust.
Use a thin frame in the plane of the discs.

Do not decrease the parasitic drag too much.
A lower $c$ increases the limit only as $c^{-1/4}$, and it decreases the agility in dives.
Keep a thrust reserve on each motor at top speed.
The top motor gets to its maximum thrust first, because it supplies thrust for speed and for
attitude control.

\section{Agility: the dive at 35\,m/s}
\label{sec:agility}

\subsection{Why the mid speed is the worst speed}

A quadrotor changes its direction when it tilts.
It tilts when its motors apply different thrusts.
The thrust of the weakest motor cannot be less than zero.
Thus, the available torque increases with the total thrust~\cite{mahony2012}.

In a dive, gravity and drag both act along the flight path.
Thus, the motors supply only a small total thrust.
At $35\ms$ in a $10^\circ$ dive, the minimum total thrust is approximately $7$\,m/s$^2$.
This is less than the thrust necessary to hover.

At the same speed, a change of target acceleration turns the thrust axis fastest.
At $55\ms$, the drag is approximately $24$\,m/s$^2$.
The total thrust is then large, and the attitude changes are small.
\Cref{tab:dive} shows the result.

\begin{table}[h]
\centering
\caption{Torque per motor in a $10^\circ$ dive. Reference quadrotor, jerk $5$.}
\label{tab:dive}
\begin{tabular}{@{}lccc@{}}
\toprule
 & $20\ms$ & $35\ms$ & $55\ms$ \\
\midrule
Available & $2.8$\,N & $\mathbf{2.4}$\,N & $7.6$\,N \\
Necessary & $1.5$\,N & $\mathbf{6.1}$\,N & $1.5$\,N \\
\bottomrule
\end{tabular}
\end{table}

\subsection{Effect of each parameter}

\Cref{tab:agility} shows that yaw torque is the limit.
Roll torque and pitch torque are not the limit.
Our method locks the frame to the side direction of the target.
Thus, a change of thrust axis causes a yaw acceleration that increases with the bank angle.

Two changes correct this problem.
The first change is a larger yaw reaction coefficient.
Use propellers and motors with more reaction torque per newton, or tilt the motors.
The second change is a lower yaw inertia.
Put the mass near the centre.

A longer arm helps only roll and pitch.
More thrust does not help, because the problem is the low minimum thrust.

\begin{table}[h]
\centering
\caption{Certified $v_{\max}$ in m/s. Dives to $10^\circ$, jerk $5$.}
\label{tab:agility}
\begin{tabular}{@{}lll@{}}
\toprule
Parameter & Values & Certified \\
\midrule
Yaw coefficient (m) & $0.015 \,/\, 0.02 \,/\, 0.03$ & $30.9 \,/\, 33.4 \,/\, \mathbf{52.1}$ \\
Inertia factor & $0.7 \,/\, 1.0 \,/\, 1.4$ & $\mathbf{52.1} \,/\, 33.4 \,/\, 30.1$ \\
Arm length (m) & $0.15 \,/\, 0.18 \,/\, 0.22$ & $33.2 \,/\, 33.4 \,/\, 33.5$ \\
Thrust-to-weight ratio & $4$ to $8$ & $33.4$ for each value \\
\bottomrule
\end{tabular}
\end{table}

\subsection{Less drag gives worse dives}
\label{sec:paradox}

An airframe with less drag dives worse.
At $c = 0.006$, the certified speed for dives to $5^\circ$ at jerk $5$ decreases from $52.1\ms$ to
$39.3\ms$.
The cause is the lower total thrust in a dive.
Drag is the air brake of the quadrotor.
In a pursuit, this air brake keeps the torque available.
A smaller rotor-disc area has the same effect (\cref{sec:checklist}).

\subsection{Wings make the problem worse}

We made a prototype of a quadrotor with wings.
Its arms had a cambered section that gives lift in the direction of the camera.
The lift decreases the load on the motors.

In a dive at $35\ms$, the minimum motor thrust decreased to approximately zero.
The available torque decreased to $0.01$\,N to $0.28$\,N, but the necessary torque was $0.52$\,N.
For all airframes, the bank angle changes by approximately $5.9^\circ$ for each m/s$^2$ of side
acceleration.
Wings help only with control surfaces, and with these surfaces the vehicle is not a quadrotor.

\section{Method: copy the target path from below}
\label{sec:strategy}

Our method has four parts.
The planner uses a three-tap filter that is exact for constant acceleration.
This filter gets the target positions with a delay of four frames at $240$\,Hz.
The quadrotor flies the filtered target path, at a fixed distance straight below it.
The thrust axis follows the planned acceleration, and the yaw follows the side direction of
the target.
The hinge points the camera straight up, in the limits of its range.

We tested five other methods.
\Cref{tab:alternatives} gives the results.
No method was better than our method.

\begin{table}[h]
\centering
\caption{Other pursuit methods and their results.}
\label{tab:alternatives}
\begin{tabularx}{\linewidth}{@{}lY@{}}
\toprule
Method & Result \\
\midrule
Linear filters with lag correction & With a delay of four frames, the filter must predict. This makes the filter gain $\sum|h_i|$ approximately 2 or more. Thus the planned jerk and snap stay larger than the jerk and snap of the target. \\
Nonlinear planner with a jerk limit & Against bad targets, the lag was 39\,m to 325\,m and $13\ms$ to $35\ms$. The target turns back when the planner is almost at the target. \\
Free yaw, mean or larger formations & These methods decrease the torque, but they increase the bank angle, the thrust or the lag. \\
Weave to decrease rotor-disc drag & For this weave, the motors must supply thrusts from $-34$\,N to $+78$\,N (\cref{sec:speed}). \\
Wings & The motors stop in dives (\cref{sec:agility}). \\
\bottomrule
\end{tabularx}
\end{table}

The proofs also show a negative result.
No proof can force a torque of practical size on the quadrotor.
The detection range is 570\,m and the lens is $120^\circ$.
With these values, the quadrotor can respond in 7\,s to 30\,s.
This forces a body rate of only $0.002$ to $0.007$\,rad/s.

Thus, the agility limits come from our tight formation, not from physics.
A looser formation can use range and field of view as a margin.
But we did not find a looser formation that gives better results.

Put your money and effort into hardware margin and into a clean planner at a high rate.
Do not put them into complex path design~\cite{mellinger2011}.

\section{Rules for operation}

Use these rules in operation:
\begin{enumerate}
\item Fly below and behind the target, at a distance of 60\,m to 100\,m.
Get into this position before the target gets to its top speed.
\item Expect the most difficult condition at mid speed in a dive.
A target that dives at $30\ms$ to $40\ms$ and starts a turn uses the most margin.
\item Know the two target movements that no quadrotor can follow.
The first movement is level flight above the speed limit, which is $57.7\ms$ for the reference
quadrotor.
The second movement, in the isotropic model, is a vertical descent faster than $\sqrt{g/c}$.
This speed is $35\ms$ at $c = 0.008$.
\item Do not follow from above at a height of more than $D\sin e^*(v)$.
If you do, the quadrotor goes past the target.
\item Treat the planner rate as a design value.
The proved bounds use $120$\,Hz to $480$\,Hz with a delay of four frames.
A slower rate increases the change between frames.
\end{enumerate}

\section{Design checklist}
\label{sec:checklist}

\Cref{tab:checklist} gives the design values for a target at $55\ms$.
We tested the recommended set of changes as one unit (\cref{tab:package}).
This set has a thrust-to-weight ratio of 7, $k_n = 0.005$, a yaw coefficient of $0.03$ and
$30\%$ less inertia.

The parameters have an effect on each other.
Without less inertia, the certified speed for dives to $10^\circ$ at jerk $5$ decreases to $35.5\ms$.
The cause is the smaller $k_n$, which also decreases the total thrust in a dive.
Thus, always test the full set of changes, not each change alone.
Above approximately $59\ms$ in $10^\circ$ dives, the tilt limit of $85^\circ$ becomes the limit.

\begin{table}[h]
\centering
\caption{Recommended set of changes, tested with the proved bounds. Values in m/s.}
\label{tab:package}
\begin{tabular}{@{}lcc@{}}
\toprule
 & Dives to $5^\circ$ & Dives to $10^\circ$ \\
\midrule
Proved limit & $73.9$ & $73.9$ \\
Certified, jerk to $3$ & $68.2$ & $58.6$ \\
Certified, jerk to $5$ & $68.1$ & $58.4$ \\
\bottomrule
\end{tabular}
\end{table}

\begin{table}[h]
\centering
\caption{Design checklist for a target at $55\ms$.}
\label{tab:checklist}
\small
\begin{tabularx}{\linewidth}{@{}rlYY@{}}
\toprule
 & Item & Design value & Reason \\
\midrule
1 & Thrust-to-weight ratio & 7 or more (6 or more with low $k_n$) & Approximately $5\ms$ for each unit \\
2 & Rotor-plane area & $k_n$ not more than $0.005\,\mathrm{m}^{-1}$ & Limit increases as $k_n^{-1/4}$ \\
3 & Yaw authority & coefficient $0.03$\,m or more, or tilted motors & Steep dives from $33$ to $52\ms$ \\
4 & Yaw inertia & not more than $0.014$\,kg\,m$^2$ at 1.5\,kg & Same effect as item 3 \\
5 & Camera & top, rear; hinge $30^\circ$ to $70^\circ$, or fixed at $50^\circ$ with a $120^\circ$ lens & The camera is then not a limit \\
6 & Planner & $240$\,Hz, delay of four frames, exact three-tap filter & Certified configuration \\
7 & Motor reserve & $1$ to $2$\,m/s$^2$ of total thrust at top speed & The top motor gets to its limit first \\
8 & Parasitic drag & $c$ not less than approximately $0.008\,\mathrm{m}^{-1}$ & Drag keeps the torque in dives \\
9 & Wings & none & Wings stop the motors in dives \\
10 & Tilt limit & $85^\circ$ or more & Limit above $59\ms$ in steep dives \\
\bottomrule
\end{tabularx}
\end{table}

\section{What to measure before you use these numbers}
\label{sec:measure}

\Cref{tab:measure} gives the quantities to measure on your quadrotor.
Each quantity has a large effect on the results.

\begin{table}[h]
\centering
\caption{Quantities to measure.}
\label{tab:measure}
\small
\begin{tabularx}{\linewidth}{@{}lYY@{}}
\toprule
Quantity & Method & Reason \\
\midrule
$k_n$ and $c$ & Wind tunnel or coast-down test at different pitch angles. Measure the drag along and normal to the rotor plane. & These values set the limit to the power $1/4$. \\
Thrust and airspeed & Thrust stand in a wind tunnel & The model keeps the maximum motor thrust constant. At $50\ms$, the airflow changes it. \\
Yaw reaction coefficient & Torque stand & This value sets the dive performance. \\
Inertia & Bifilar pendulum & Same reason \\
Target jerk and snap & Flight records of the target type & Each certified speed is for one jerk value. \\
\bottomrule
\end{tabularx}
\end{table}

\section{Conclusion}

Hardware margin, not the pursuit method, limits a quadrotor that follows a fast fixed-wing
aircraft.
The thrust-to-weight ratio and the rotor-plane area set the top speed.
The proved speed limit is less than $1\ms$ from steady flight.
Yaw authority and yaw inertia set the agility in a dive at mid speed.
The camera position removes the camera as a limit.
A companion paper gives the formal proofs.

\section*{Reproducibility}

The repository is \url{https://github.com/AbdullahAlmanei/vision-tracking-lean}.
Use these commands to get the results again:
\begin{itemize}\raggedright
\item Proofs: \texttt{lake build}, then \texttt{scripts/axcheck.sh master\_sufficient3 p3\_bracket\_power}.
\item Parameter tables: \texttt{scripts/paper\_sweeps.py}.
\item Bounds: \texttt{scripts/gap\_v3.py} and \texttt{scripts/gap\_v3.py kn}.
\item Weave study: \texttt{scripts/explore/weave.py}.
\end{itemize}

\begin{thebibliography}{9}\small
\bibitem{lean4} L.~de~Moura and S.~Ullrich. The Lean 4 theorem prover and programming language. In
\emph{CADE-28}, 2021.
\bibitem{mathlib} The mathlib Community. The Lean mathematical library. In \emph{CPP}, 2020.
\bibitem{hoffmann2007} G.~M.~Hoffmann, H.~Huang, S.~L.~Waslander and C.~J.~Tomlin. Quadrotor helicopter
flight dynamics and control: theory and experiment. In \emph{AIAA Guidance, Navigation and Control
Conference}, 2007.
\bibitem{faessler2018} M.~Faessler, A.~Franchi and D.~Scaramuzza. Differential flatness of quadrotor
dynamics subject to rotor drag for accurate tracking of high-speed trajectories. \emph{IEEE Robotics and
Automation Letters}, 3(2), 2018.
\bibitem{mahony2012} R.~Mahony, V.~Kumar and P.~Corke. Multirotor aerial vehicles: modeling, estimation,
and control of quadrotor. \emph{IEEE Robotics \& Automation Magazine}, 19(3), 2012.
\bibitem{mellinger2011} D.~Mellinger and V.~Kumar. Minimum snap trajectory generation and control for
quadrotors. In \emph{IEEE ICRA}, 2011.
\end{thebibliography}

\end{document}
